The experimental campaign can be read as a sequence of falsifiable claims. Claim~A is that a PID-held-out prefix model can exceed published accuracy references while remaining calibrated on NYC Open Data. Claim~B is that budgeted VoI inspection improves detection metrics relative to fixed ranking under equal $B$. Claim~C is that disagreement priors change inspection outcomes relative to consensus priors. The TEST evidence supports Claim~A, rejects Claim~B under the clamped evidence model, and fails to support Claim~C at the alert-metric level. That pattern is more informative than a uniformly positive result set because it separates what the data can identify from what the method can merely compute~\cite{Song2021,Memarzadeh2016,Egwim2021}.

Table~\ref{tab:champion} should be read as a metric pack, not as accuracy alone. Balanced accuracy, precision, recall, specificity, F1/F2, ROC-AUC, PR-AUC, Brier, log-loss, ECE, and Cohen $\kappa$ jointly show whether the model is a trivial majority classifier or a calibrated triage score~\cite{Saito2015,Gneiting2007,Cohen1960,Hand2009}.

The confusion matrix in Fig.~\ref{fig:cm} makes the error trade visible in counts rather than ratios. With $16$ false positives and $17$ false negatives at $\tau{=}0.51$, the operating point is nearly symmetric in absolute errors despite a $4{:}1$ non-severe to severe ratio. That symmetry is a consequence of VAL tuning with a soft recall floor, not an accident of TEST peeking. If an agency preferred fewer false alarms, $\tau$ would rise and net benefit in Fig.~\ref{fig:dca} would indicate whether the preference remains advantageous~\cite{Vickers2006}.

Reliability diagrams such as Fig.~\ref{fig:cal} diagnose local miscalibration that aggregate Brier scores can average away. The observed mild over-confidence in a high-probability bin is consistent with tree-ensemble behavior on tabular features and does not overturn the low overall ECE~\cite{Murphy1973,Gneiting2007}. For deployment, the practical response is monitoring: if future cohorts drift, recalibration on a rolling VAL window can be introduced without changing the PID split discipline.

Budgeted recall of $0.450$ at high precision looks weak beside the champion's $0.845$ recall until the tasks are distinguished. The BN scoreboard evaluates alerts after at most $B$ factor inspections with absence-clamped evidence, not after consuming the full prefix feature vector. Many severe projects simply lack inspectable reason diversity, so the BN cannot conjure evidence that the administrative record never offered. The champion exploits cumulative variance trajectories that do not require naming a factor~\cite{Luu2009,Tagharobi2025}.

Mean inspections-to-first-detection near $1$ among alerting severe projects indicates that, when an alert occurs under budget $B{=}4$, it usually occurs after the first informative inspection. Increasing $B$ cannot create decisive evidence if the inspectable set is empty or uninformative, which is why Table~\ref{tab:fullB} is flat in classification metrics across budgets~\cite{Fauriat2020,Yuan2021}.

Expected utility and realized loss in Fig.~\ref{fig:util} close the loop between VoI's internal objective and external evaluation. Because the same asymmetric costs enter both, a policy that maximized myopic VoI but performed poorly on realized loss would signal an implementation bug or a mismatch between myopic and multi-step value. The observed tie across policies shows that all policies face the same sparse evidence menu~\cite{Howard1966,Memarzadeh2016}.

Literature on construction AI adoption cautions that algorithmic novelty often outpaces evaluation rigor~\cite{Darko2020,PanYZhang2021,Abioye2021,Regona2024}. The present manuscript responds by refusing to treat VoI as automatically beneficial and by publishing the policy tie. A negative comparative result under a clearly described evidence model is preferable to a positive result obtained by relaxing leakage controls or by tuning on TEST.

From the standpoint of mathematical provenance, \eqref{eq:noisyor}--\eqref{eq:noisyor2} are adopted Noisy-OR identities~\cite{Heckerman1996,Onisko2001}; \eqref{eq:voi}--\eqref{eq:voi2} are adapted VoI identities with an alert-loss action space~\cite{Howard1966,Song2021,Memarzadeh2016}; the PID split protocol and metric pack are adapted from predictive-study practice in delay ML and probabilistic forecasting~\cite{Egwim2021,Tagharobi2025,Gneiting2007,Vickers2006}. The derived contribution is the coupling and the experimental measurement of policy identifiability on open capital-project episodes.

Parameter provenance for priors follows empirical TRAIN frequencies rather than informal expert interviews. Interview-based CPTs can encode institutional knowledge that open data omit~\cite{Khodabakhshian2025,Luu2009}, but they also introduce elicitation noise and make multi-city replication harder. Empirical priors make replication straightforward at the cost of depending on historical reason-writing practices.

External validity remains the largest open scientific risk. NYC capital-project administration, funding instruments, and reason taxonomies differ from marine construction~\cite{Comu2023}, mass housing~\cite{Khodabakhshian2025}, and UK progress portfolios~\cite{Tagharobi2025}. Optimism-bias literature further warns that institutional forecasting cultures differ~\cite{Flyvbjerg2009,Love2015}. The numeric operating points are therefore local facts until replicated.

Baseline fairness must be stated carefully. The BN baselines (fixed ranking and static BN) are implemented inside the same evidence-clamping and thresholding machinery as VoI, which is the correct fairness setup for RQ1--RQ2. The early-warning literature comparison is a scoped numeric comparison to published accuracies, not a same-codebase bake-off~\cite{Khodabakhshian2025,Tagharobi2025,Egwim2021}.

Sensitivity to early fraction (Table~\ref{tab:early}) shows that discrimination remains strong even when the prefix is shortened to $0.40$ of updates, with accuracy $0.902$ and AUC $0.953$. Agencies that value earlier intervention can trade a few points of accuracy for more remaining updates~\cite{Tagharobi2025}.

The bootstrap design uses $400$ resamples of TEST PIDs, which is enough for stable percentile intervals on accuracy and AUC given $n{=}550$. It is not a hierarchical bootstrap over agencies or over time blocks. The intervals answer how uncertain the headline TEST metrics are to resampling of the same TEST population~\cite{GneitingRaftery2007,Hand2009}.

Integrity auditing during ingest inspected duplicates, impossible negatives, and missing agency fields under scripted rules; blank reasons were dropped from mapping rather than imputed. The audit does not prove that every administrative reason is a faithful causal statement. It proves that the modeling sample was not silently hand-edited to improve metrics~\cite{Woldesenbet2016}.

Administrative delay reasons in capital dashboards are often short, repetitive, and written for compliance rather than for causal analysis. That writing culture compresses the inspectable factor set and thereby compresses VoI's opportunity set. The empirical tie in Table~\ref{tab:fullB} is therefore partly a data-regime result~\cite{Woldesenbet2016,Assaf2006}.

Public agencies face reputational and fiscal costs when severe delays are missed, which motivates $C_{\mathrm{MISS}}>C_{\mathrm{FA}}$. The specific $5{:}1$ ratio is a transparent placeholder for that asymmetry. Decision-curve analysis then allows stakeholders to explore other preferences without rerunning the entire BN whenever only $\tau$ changes~\cite{Vickers2006,Howard1966,Love2015}.

Approximate BN inference introduces an extra error term that can interact with policy comparisons. With $256$ configurations, that error term is avoidable. Exact inference therefore strengthens the negative VoI-comparative result: the tie is not an artifact of loopy belief propagation~\cite{Heckerman1996,Onisko2001}.

Row-wise random splits on multi-period project data can place early and late updates from the same PID into both TRAIN and TEST, inflating accuracy. PID-level splitting is the primary leakage control in this study and should be treated as a methodological requirement for similar open schedule datasets~\cite{Tagharobi2025,Egwim2021}.

Defining severe delay as a top quartile constructs a $0.20$ base rate by design among eligible PIDs. That definition is transparent and reproducible, but it is not identical to contractual liquidated-damages events. Metric interpretation must respect that label semantics~\cite{Saito2015,Hand2009,Flyvbjerg2009}.

Eight factors cannot encode every contractual nuance. They can encode a stable interface between text and BN parents. Expanding the ontology without expanding inspectable evidence would likely worsen sparsity and further reduce VoI differentiation~\cite{Onisko2001,Luu2009}.

A low Brier score indicates that predicted probabilities match outcomes in mean squared error. Operators can therefore treat $p$ as a triage score rather than as a mere ranking feature, provided cohort shift is monitored~\cite{Brier1950,Murphy1973,Gneiting2007}.

ROC-AUC ignores the cost structure of thresholds. Net benefit does not. Fig.~\ref{fig:dca} is therefore not redundant with Fig.~\ref{fig:rocpr}; it answers a different decision question~\cite{Vickers2006,Hand2009}.

Even when disagreement VoI ties consensus VoI, storing $\pi^{\mathrm{d}}$ beside $\pi^{\mathrm{c}}$ documents where the model is least certain a priori. That metadata can guide future elicitation campaigns even if it does not change today's alerts~\cite{Khodabakhshian2025,Memarzadeh2016}.

Schedule variance itself is a noisy proxy for managerial distress. Positive cumulative variance can reflect re-baselining, weather, funding gaps, or reporting artifacts. The champion treats variance trajectories as predictive signals without asserting that every variance unit is causal~\cite{Gneiting2007,Kim2016,Nasir2003}.

Bayesian delay models sometimes emphasize structure learning. Structure learning is powerful when expert graphs are unavailable, but it can be unstable at small $N$~\cite{Zhang2021BNMCS,Khodabakhshian2025}. DRIN fixes a compact Noisy-OR structure adapted from survey taxonomies and concentrates experimental power on inspection policy and early warning.

Marine and tunnel BBN applications show that domain physics and safety events can enrich evidence beyond administrative text~\cite{Comu2023,Wu2015,Wang2016,Leu2012}. Those richer evidence channels are exactly where VoI differentiation is more likely to appear. The NYC results bound VoI value in reason-sparse capital dashboards; they do not deny VoI value in sensor-rich domains.

Fuzzy and neuro-fuzzy risk models remain relevant when linguistic expert assessments dominate~\cite{Ebrat2014,Islam2017}. They are less natural when the primary evidence is coded open data with binary factor presence. DRIN's probabilistic stack is chosen for the latter regime, while acknowledging hybrid fuzzy-BN systems in elicitation-heavy settings~\cite{Qazi2016,Zou2017}.

Cohen's $\kappa{=}0.812$ indicates strong chance-corrected agreement at the chosen threshold~\cite{Cohen1960}. Matthews' correlation is related in spirit for binary tasks~\cite{Matthews1975}. These coefficients are secondary to proper scoring for probability quality but help communicate classifier agreement to audiences accustomed to categorical reliability metrics.

The index of experiments E1--E6 is deliberately smaller than a mechanical E1--E8 checklist. Scalability beyond eight factors is analytically bounded in the text, and multi-city external validation is limited by available authoritative open endpoints under the same schema. Omitting a weak synthetic second dataset is preferable to manufacturing an external-validation theater~\cite{Woldesenbet2016}.

Hyperparameter search was intentionally light. The dominant generalization risk in this study is leakage, not underfitting of tree depth. The reported champion reflects that priority ordering~\cite{Egwim2021,Friedman2001}.

When policies tie, one might ask whether VoI should be removed from the system. If the sole goal is TEST recall at $B{\leq}4$ on NYC reasons, fixed ranking is enough. If the goal is an auditable decision-analytic layer that can exploit denser evidence in future datasets without rewriting the stack, retaining VoI remains rational~\cite{Howard1966,Song2021}.

Open administrative data are abundant relative to expert-only case studies, yet semantically shallow relative to site instrumentation. Models that ignore shallowness overclaim; models that ignore abundance underuse public evidence. DRIN's clamped-evidence BN and prefix-based champion are two responses tuned to that trade-off~\cite{Woldesenbet2016,Luu2009,Egwim2021,Regona2024}.

A practical deployment reading follows. Portfolio managers can sort projects by calibrated $P(Y{=}1)$ from prefix features and allocate scarce senior attention to the upper tail. Analysts assigned to those projects can then use the BN to record which factor hypotheses were inspected and how the posterior moved. That workflow does not require VoI to beat fixed ranking; it requires probabilities and factor posteriors to be coherent~\cite{Vickers2006,Luu2009}.

From a theoretical perspective, the flat scoreboard is consistent with a near-zero VoI gap whenever the feasible observation set is small and the alert loss is locally constant in $p$. In that regime, expected loss reduction from reordering inspections vanishes even though the VoI formula remains well defined~\cite{Howard1966,Song2021,Malings2016}. The mathematics is not wrong; the opportunity set is thin.

Lead-time results refine that deployment reading. Caught severe projects retain on average nearly three remaining updates after the early window, which is actionable for many capital programs but not a guarantee of recovery. Missed severe projects show that some severe trajectories remain quiet in early variance space, so prefix-only warning has a hard envelope~\cite{Tagharobi2025,Egwim2021}.

Mapping-quality results refine the BN reading. NLP-only macro-F1 against regex weak labels is too low to serve as a sole extractor, whereas regex-first recovery on informative rows is stable. The BN therefore inherits an auditable evidence layer rather than an unconstrained text classifier~\cite{Onisko2001,Woldesenbet2016}.

Comparison discipline against published accuracies remains essential. The $0.85$ and $0.89$ references come from elicitation-augmented BN risk modeling and stage-progress prediction, respectively~\cite{Khodabakhshian2025,Tagharobi2025}. They are yardsticks, not identical-task replicates. Bootstrap intervals keep the yardstick comparison from collapsing to a single fragile point.

The manuscript's quantitative figures each support a distinct inferential job: scorecard overview, literature yardstick, metric pack, ranking curves, error counts, calibration, net benefit, uncertainty, comparison table, budget curves, detection steps, lead time, mapping quality, and utility. Removing any one would remove a corresponding claim pathway~\cite{Gneiting2007,Vickers2006,Saito2015}.

Finally, the research gap stated in Section~\ref{sec:lit} is only partially closed by one metropolitan study. What is closed is the absence of a joint, leakage-safe experimental account of calibrated early warning and budgeted BN inspection on the same open capital-project PIDs. What remains open is whether the VoI-policy tie persists in denser evidence regimes and in other cities. That residual openness is a boundary of the evidence, not a failure of the experimental design~\cite{Memarzadeh2016,Regona2024,Abioye2021}.

Causal surveys consistently show multi-factor delay generation~\cite{Assaf2006,Sambasivan2007,Doloi2012}, which justifies a multi-parent BN child. They do not justify assuming that every factor is inspectable at every decision epoch. The inspection-budget constraint is therefore not an artificial academic restriction; it mirrors staffing limits that agencies already face when investigating variance explanations.

ERIC-S and related schedule-risk systems demonstrated that probabilistic risk registers can be operationalized~\cite{Nasir2003}. Bayesian duration updating showed that beliefs about remaining time can be revised as evidence arrives~\cite{Kim2016}. DRIN inherits those intuitions but changes the decision variable from duration percentiles alone to a budgeted factor-inspection policy tied to severe-delay alerts.

Strictly proper scoring is central to the early-warning claims. A model could beat literature accuracy while emitting unusable probabilities. Brier and log-loss guard against that failure mode, and ECE diagnoses bin-wise reliability~\cite{Brier1950,GneitingRaftery2007,Gneiting2007}. Without those metrics, RQ3 would be only a discrimination question.

Decision-curve net benefit further guards against threshold naivety. An agency that cannot tolerate false alarms needs a different $\tau$ than an agency that cannot tolerate missed severe delay. Fig.~\ref{fig:dca} shows that the champion retains value across a band, which is stronger operational evidence than a single accuracy number~\cite{Vickers2006}.

The disagreement prior experiment is intentionally modest. It does not claim a full multi-expert elicitation protocol. It asks whether a transparent widening of uncertain base rates changes held-out alert metrics under VoI. On NYC TEST data, it does not, which is a useful negative for teams considering expensive elicitation solely to feed myopic VoI~\cite{Khodabakhshian2025,Song2021}.

Computationally, the eight-factor design keeps posterior and VoI queries interactive. That property matters for adoption more than asymptotic complexity arguments. If future ontologies grow, the same experimental template can swap exact enumeration for factored inference while preserving the scoreboard protocol~\cite{Heckerman1996,Onisko2001}.

The absence of a second primary dataset is a limitation, not an evasion. Mixing UK GMPP summaries into TEST would create an illusory external validation with mismatched schemas and labels. Keeping those summaries out preserves the integrity of the PID-held-out NYC claims~\cite{Woldesenbet2016,Tagharobi2025}.

Reproducibility requires more than sharing accuracy. The released scoreboard CSV, bootstrap summary, and figure-generation script allow independent recomputation of Tables~\ref{tab:b2}--\ref{tab:fullB} and of the literature-yardstick figures. That evidence trail is part of the contribution insofar as it makes the negative VoI result auditable.

In closing the Results interpretation, the strongest supported finding is calibrated early warning under PID-held-out open data; the strongest supported boundary is VoI-policy non-identifiability under clamped sparse reasons. Both belong in the abstract and conclusion because both are experimentally earned~\cite{Egwim2021,Memarzadeh2016,Gneiting2007}.

Survey-derived cause lists remain indispensable for ontology design, but they rarely specify a sequential inspection objective. Assaf and Al-Hejji, Sambasivan and Soon, and Doloi et al.\ establish which factors recur; they do not establish which factor should be verified first under a numeric budget~\cite{Assaf2006,Sambasivan2007,Doloi2012}. Bridging that gap requires an explicit decision model, which is why VoI enters the formulation even when the eventual empirical comparison is negative.

Infrastructure overrun research adds a behavioral layer. Flyvbjerg's analyses of optimism bias and Love et al.'s synthesis of transport overruns imply that forecast error is structured by incentives and institutions~\cite{Flyvbjerg2009,Love2015}. Open schedule dashboards capture some of that structure through variance trajectories, but they under-capture incentive narratives. Early-warning models can still exploit the captured trajectories; BN factor models should remain modest about what reason text can identify.

In reliability VoI, information value is often largest when prior uncertainty is high and observations are decisive~\cite{Memarzadeh2016,Malings2016,Yuan2021}. NYC reason inspections are frequently neither: many PIDs have few inspectable factors, and confirming an already-suspected factor may not move $p$ across $\tau$. The scoreboard's flatness is therefore consistent with VoI theory rather than contradictory to it~\cite{Song2021,Fauriat2020}.

Calibration theory clarifies a different success mode. Sharp, calibrated probabilities create decision value even when a secondary inspection policy adds little~\cite{Gneiting2007,GneitingRaftery2007,Brier1950}. That is the empirical pattern here: the champion supplies decision value; budgeted BN policies supply auditability with limited differential detection gain.

The early-fraction sweep connects to stage-wise progress prediction literature, where information content increases as projects advance~\cite{Tagharobi2025}. Maintaining AUC above $0.95$ at early fraction $0.40$ indicates that useful signal appears relatively early in many NYC trajectories, which is encouraging for triage and consistent with the lead-time distribution for caught cases.

Imbalance-aware evaluation is mandatory at base rate $0.20$. Accuracy alone can look strong for a weak model; PR-AUC $0.944$ and balanced accuracy $0.905$ show that severe-class ranking and balanced correctness are also strong~\cite{Saito2015,Hand2009,Egwim2021}. The confusion counts then confirm that the threshold does not hide a one-sided error pattern.

For BN parents, regex-first mapping is an adopted engineering pattern rather than a learned causal discovery procedure. It trades recall of rare linguistic variants for precision and auditability. In Noisy-OR settings with limited data, such constrained evidence models are often preferable to unstable high-capacity extractors~\cite{Onisko2001,Heckerman1996}.

Disagreement widening can be viewed as a sensitivity analysis embedded in the prior. If alerts change under $\pi^{\mathrm{d}}$, prior uncertainty is decision-relevant; if they do not, prior uncertainty is dominated by likelihood evidence or by clamping. TEST supports the latter reading for alert metrics, while still leaving room for disagreement to matter in posterior narratives shown to analysts~\cite{Khodabakhshian2025}.

The scoped literature beat should be communicated with the same caution in talks and abstracts as in the Results text. Saying that the champion exceeds published accuracy references under a PID-held-out open-data protocol with calibration diagnostics is supported. Saying that it outperforms all prior delay models is not~\cite{Khodabakhshian2025,Tagharobi2025,Egwim2021}.

Agency heterogeneity is an unmodeled factor. Different NYC agencies may write reasons differently and may experience different variance regimes. The present study pools agencies after including agency indicators in the champion features, but it does not claim agency-invariant BN CPTs. Stratified agency experiments are a natural extension when sample sizes per agency permit stable evaluation~\cite{Woldesenbet2016}.

Temporal non-stationarity is another unmodeled factor. Capital programs evolve with mayoral priorities, funding shocks, and construction-market conditions. A single random PID split estimates contemporaneous generalization, not forward-in-time generalization. Temporal holdout would be a stronger external-validity test and is recommended for follow-on work~\cite{Tagharobi2025,Love2015}.

The utility of exact enumeration extends beyond speed. It makes educational and auditing use cases possible: an analyst can in principle recompute a posterior by listing compatible configurations. That transparency is difficult to preserve with black-box sequence models, which is one reason the BN layer remains valuable beside the champion~\cite{Luu2009,Heckerman1996}.

Related hybrid risk research using fuzzy inference encodes linguistic uncertainty that open-data binaries discard~\cite{Ebrat2014,Islam2017,Qazi2016}. If future datasets include structured expert assessments alongside schedule updates, a hybrid fuzzy-BN or fuzzy-VoI extension could be justified. With the current endpoints, introducing fuzzy membership without new evidence would be cosmetic.

BIM-oriented risk reviews emphasize information integration across design and construction~\cite{Zou2017}. DRIN integrates only schedule-history and budget covariates because those are the authoritative open endpoints used here. Richer BIM-linked features could raise early-warning lead time and expand inspectable evidence, potentially reopening VoI differentiation.

Safety-oriented construction BBNs for tunnels and equipment accidents demonstrate evidence updating under field observations that are more decisive than administrative reasons~\cite{Wu2015,Wang2016,Leu2012}. Those papers indirectly explain why VoI shines in some construction subdomains and stalls in others: observation quality dominates policy sophistication.

Scientometric reviews chart rapid AI uptake and uneven rigor~\cite{Darko2020,PanYZhang2021,Abioye2021,Regona2024}. The evaluation pack in this paper---PID split, proper scoring, decision curves, bootstrap CIs, and a full budgeted scoreboard---is a direct response to that unevenness. Method sections that only name algorithms without such a pack leave RQ-level claims under-identified.

Friedman-style gradient boosting provides the statistical backbone of the champion~\cite{Friedman2001}. The hist-based implementation is an engineering detail; the scientific point is that a strong tabular ensemble plus leakage-safe features can yield calibrated early warning without deep sequence modeling on this dataset.

Howard's information-value framing remains the correct conceptual parent for the inspection layer even after a negative comparative result~\cite{Howard1966}. A theory can be correct and still yield near-ties under thin opportunity sets. Reporting those ties is how applied VoI research stays honest~\cite{Memarzadeh2016,Song2021}.

The conclusion's restraint follows from these interpretations. The study does not claim a universal inspection policy upgrade. It claims a measured operating envelope on NYC Open Data and a reproducible experimental template for testing whether that envelope moves when evidence becomes denser or when cities change~\cite{Woldesenbet2016,Egwim2021,Luu2009}.
